% !TeX root = ../QFT_main.tex
\documentclass[../QFT_main.tex]{subfiles}
\begin{document}
%=========================================================
%=========================================================
\chapter{Symmetries}\label{ch:symmetries}
%========================================================
%=========================================================





%--------------------------------------------------------
%=========================================================
\section{Discrtete Symmetries}
%=========================================================
%------------------------------------------------------------



%--------------------------------------------------------
%=========================================================
\section{Continuous Symmetries}
%=========================================================
%----------------------------------------------------------



Noether's theorem is a fundamental result in theoretical physics that establishes a deep connection between symmetries and conservation laws.
It states that for every continuous symmetry of the action of a physical system, there exists a corresponding conserved quantity. In other words, if the laws of physics remain invariant under a continuous transformation (such as translations, rotations, or gauge transformations), then there is an associated physical quantity that remains constant over time.
We first discuss the classical version.

\begin{mytheorem}[Classical Noether's 1st theorem \questioninline{polish}]
Consider the action for a classical field theory in $d$ spacetime dimensions,
{\small\begin{align}
    S[\phi] = \int_\Omega \mathcal{L}(\phi^a, \partial_\mu \phi^a,x^\mu)\,\,d^d\!x\,,
\end{align}}
where $\phi^a$ is any field with arbitrary spin and internal indices labelled by $a$.

Consider a continuous transformation spacetime coordinates and fields, parametrized infinitesimally as
{\small\begin{align}
    x^\mu \mapsto \tilde{x}_\epsilon^\mu\simeq  x^\mu +\epsilon \delta x^\mu\,,
    \qquad
    \phi^a(x)\mapsto \tilde{\phi}_\epsilon^a(x)\simeq \phi^a(x)+\epsilon \Delta \phi^a(x)\,\,,
\end{align}} 
where $\epsilon \Delta \phi^a(x)\simeq \tilde{\phi}_\epsilon^a(x)-\phi^a(x)$ is the variation of the field at \emph{fixed} spacetime point $x$.
More rigorously, it is simply the derivative wrt the parameter $\epsilon$ of the transformed function $\Delta f(x):=\lim_{\epsilon\to0}\frac{\tilde{f}_\epsilon(x)-f(x)}{\epsilon}$.
Note that variations at fixed spacetime point and derivatives commute, simply because partial derivatives commute by Schwarz's theorem, namely
{\small\begin{align}\label{eq:noether_derivative_vs_variation_fixed_point_bis}
\begin{aligned}
    \Delta(\partial_\rho\phi^a)(x)&:= \lim_{\epsilon\to0} \frac{1}{\epsilon} \left[\tfrac{\partial\tilde{\phi}_\epsilon^a(y)}{\partial y^\rho}_{\mid y=x}-\tfrac{\partial{\phi}^a(y)}{\partial y^\rho}_{\mid y=x}\right]
    =
    \tfrac{\partial}{\partial\epsilon}\tfrac{\partial}{\partial y^\rho}\tilde{\phi}_\epsilon^a(y)_{\mid \substack{y=x\\ \epsilon=0}}
    =
    \tfrac{\partial}{\partial y^\rho} \tfrac{\partial}{\partial\epsilon}\tilde{\phi}_\epsilon^a(y)_{\mid \substack{y=x\\ \epsilon=0}}
    = \tfrac{\partial}{\partial x^\rho}(\Delta\phi^a)(x)\,.
\end{aligned}
\end{align}}
We also remark this is not the case for total variations $\epsilon\delta \phi^a:=\tilde{\phi}^a(\tilde{x})-\phi^a(x)$, that is $\delta(\partial_\rho\phi^a)\neq \partial_\rho \delta \phi^a$ in general.

The variations of the Lagrangian density at fixed spacetime point and under the combined transformation of spacetime coordinates and fields are given by
{\small\begin{align}\label{eq:variations_lagrangian_bis}
    \Delta\mathcal L
    &= \tfrac{\partial\mathcal{L}}{\partial \phi^a}_{|\partial\phi, x} \Delta \phi^a + \tfrac{\partial\mathcal{L}}{\partial (\partial_\mu \phi^a)}_{|\phi,x} \Delta (\partial_\mu\phi^a),
    \qquad
    \delta \mathcal{L} = \tfrac{\partial\mathcal{L}}{\partial \phi^a}_{|\partial\phi, x} \Delta \phi^a + \tfrac{\partial\mathcal{L}}{\partial (\partial_\mu \phi^a)}_{|\phi,x} \Delta (\partial_\mu\phi^a)
    +\tfrac{d\mathcal{L}}{dx^\mu}_{|\phi,\partial\phi}\, \delta x^\mu,
\end{align}}
where the subscript $|\phi$ or $|\partial\phi$ indicates that the \emph{functional} form $x\mapsto f(x)$ is kept fixed, but varying $x^\mu$ still produces terms $\propto \tfrac{\partial \mathcal{L}}{\partial f}\partial_\nu f$ due to the dependence of $f$ on $x$. 

The infinitesimal action is $d S= d^d\!x\, \mathcal{L}$.
Using \eqref{eq:variations_lagrangian_bis} and $\delta d^d\! x = (\partial_\mu \delta x^\mu) d^d\! x$, and then \eqref{eq:noether_derivative_vs_variation_fixed_point_bis}, its total variation under the combined transformation is rewritten as
{\small\begin{align}
\begin{aligned}
    \delta dS &= \big(\delta \mathcal{L}\big) \,\,d^d\!x + \mathcal{L}\, \, \delta d^d\!x 
    \\
    &= \bigg[\tfrac{\partial\mathcal{L}}{\partial \phi^a}_{|\partial\phi, x} \Delta \phi^a + \tfrac{\partial\mathcal{L}}{\partial (\partial_\mu \phi^a)}_{|\phi,x} \Delta (\partial_\mu\phi^a)
    +\tfrac{d\mathcal{L}}{dx^\mu}_{|\phi,\partial\phi}\, \delta x^\mu\bigg] d^d\!x + \mathcal{L} \, (\partial_\mu \delta x^\mu)\, d^d\!x
    \\
    &= \tfrac{d}{dx^\mu}\bigg[\tfrac{\partial\mathcal{L}}{\partial (\partial_\mu \phi^a)}_{|\phi,x}\Delta\phi^a+ \mathcal{L} \,\delta x^\mu\bigg] d^d\!x
    + \bigg[\underbrace{\tfrac{\partial\mathcal{L}}{\partial \phi^a}_{|\partial\phi, x}-\tfrac{d}{dx^\mu}\tfrac{\partial\mathcal{L}}{\partial (\partial_\mu \phi^a)}_{|\phi,x}}_{=0\,\, \text{on shell by Euler-Lagrange}} \bigg] \Delta\phi^a\,\, d^d\!x\,.
\end{aligned}
\end{align}}
Note the second term is precisely the Euler-Lagrange equations and thus vanishes on-shell.

Assume the transformation is a quasi-symmetry, meaning that the infinitesimal action is left invariant up to a total derivative $\delta\, dS \overset{!}{=} \partial_\mu F^\mu\,d^d\!x$for some function $F^\mu$.
The Noether current associated with this symmetry is then
{\small\begin{align}
    j^\mu = \frac{\partial \mathcal{L}}{\partial (\partial_\mu \phi^a)} \Delta \phi^a + \mathcal{L}\,  \delta x^\mu- F^\mu\,.
\end{align}}
Noether's current obeys the following off-shell identity, 
{\small\begin{align}
    \partial_\mu j^\mu =
    \Big[\underbrace{\tfrac{d}{dx^\mu}\tfrac{\partial\mathcal{L}}{\partial (\partial_\mu \phi^a)}_{|\phi,x}-\tfrac{\partial\mathcal{L}}{\partial \phi^a}_{|\partial\phi, x}}_{=0\,\, \text{on shell}}\Big] \Delta\phi^a\,,
\end{align}}
and in particular is conserved on-shell i.e. when the Euler-Lagrange equations are satisfied.

The associated charge is obtained by integrating the time component of the current over an arbitrary constant-time hypersurface $\Sigma_t$ and is indeed conserved in time upon assuming appropriate boundary conditions at spatial infinity:
{\small\begin{align}
    Q_t = \int_{\Sigma_t} d^{d-1}\!x\,\, j^0, \qquad \tfrac{dQ_t}{dt} = -\int_{\Sigma_t} d^{d-1}x\, \nabla_\ell j^\ell 
    =-\oint_{\partial\Sigma_t}\!\!\!\!dS\, n_{\Sigma_t}^i\,j^i= 0\,.
\end{align}}
More generally, we can build conserved charges by integrating the current over arbitrary spacelike hypersurfaces $\Sigma$ with normal vector $n^\mu$ and metric volume element $d\Sigma$, that is by integrating the Hodge dual of the current 1-form $j=j_\nu dx^\nu$ over $\Sigma$,
{\small\begin{align}
    Q_\Sigma = \int_\Sigma \!\!\!d\Sigma\, n_\mu j^\mu = \int_\Sigma \star j\,.
\end{align}}
Charges associated to homologous spacelike hypersurfaces $\Sigma_-$ and $\Sigma_+$ coincide by Gauss's theorem, upon assuming vanishing flux through any boundary at spatial infinity,
{\small\begin{align}
    Q[\Sigma_+]-Q[\Sigma_-]
    = \int_{\partial V}\star j
    =\int_{V} d\star j
    =0\,,
\end{align}}
where $V$ is the spacetime region bounded by$\Sigma_-$ and $\Sigma_+$.
%
\end{mytheorem}
%
\begin{proof}
We remark that the current features the variation of the field at fixed spacetime point {\small$\epsilon\Delta \phi^a(x):=\tilde{\phi}^a(x)-\phi^a(x)$}, not the variation at the transformed spacetime point $\epsilon\delta \phi^a=\tilde{\phi}^a(\tilde{x})-\phi^a(x)$.
For example, for a scalar field transforming only via pullback under spacetime transformations we have $\Delta \phi^a(x)=-\partial_\mu \phi^a(x)\delta x^\mu$, while $\delta \phi^a(x)=0$.

We also make explicit the variations of coordinates and fields to fix ideas.
Write the infinitesimal continuous transformations of spacetime coordinates and fields as
{\small\begin{align}
    x \mapsto \tilde{x}_\epsilon\equiv S_\epsilon(x)\simeq x+\epsilon \delta x\,,\qquad 
    \phi^a(x)\mapsto \tilde{\phi}_\epsilon^a(x)\equiv \exp(\epsilon \delta M)^a_b\phi^b(S_\epsilon^{-1}(x))\simeq (\operatorname{I}+\epsilon \delta M)^a_b \,\phi^b(x-\epsilon \delta x)\,\,.
\end{align}}
In particular, from the definitions, we have
{\small\begin{align}
    \Delta\phi^a(x):=\tilde{\phi}^a(x)-\phi^a(x) = \delta M^a_b \phi^b(x)-\tfrac{\partial \phi^a}{\partial x^\nu} \delta x^\nu\,,
    \qquad
    \delta\phi^a(x):=\tilde{\phi}^a(\tilde{x})-\phi^a(x) = \delta M^a_b \phi^b(x)\,.
\end{align}}
Due the change in the functional form of the field, the derivative of the transformed field is not simply the transformed derivative of the field, but rather
{\small\begin{align}
\begin{aligned}
   \tfrac{\partial\tilde{\phi}_\epsilon^a(y)}{\partial y^\rho}
   &=\tfrac{d}{dy^\rho}\big[\exp(\epsilon\delta M)^a_b \phi^b(S_\epsilon^{-1}(y))\big]
   \\
   &=
   \tfrac{\partial }{\partial z^\sigma}\big[\exp(\epsilon\delta M)^a_b \phi^b(z)\big]_{\mid z=S_\epsilon^{-1}(y)}\tfrac{\partial (S_{\epsilon}^{-1}(y)\big)^\sigma}{\partial y^\rho}
   \\
    &\simeq
    \exp(\epsilon\delta M)^a_b \tfrac{\partial \phi^b(z)}{\partial z^\sigma}_{\mid z=S_\epsilon^{-1}(y)}\big(\delta_\rho^\sigma-\epsilon \partial_\rho\delta y^\sigma\big)
    \\
    &\simeq 
    \tfrac{\partial \phi^a(y)}{\partial y^\rho}
    +\epsilon \delta M^a_b \tfrac{\partial \phi^b(y)}{\partial y^\rho}
    -\epsilon \tfrac{\partial \phi^a(y)}{\partial y^\sigma}\partial_\rho\delta y^\sigma
    -\epsilon \tfrac{\partial^2 \phi^a(y)}{\partial y^\rho \partial y^\nu} \delta y^\nu\,.
\end{aligned}
\end{align}}
As already shown in \eqref{eq:noether_derivative_vs_variation_fixed_point_bis} we can explictly verify that derivatives and variations at fixed spacetime point commute
{\small\begin{align}\label{eq:noether_derivative_vs_variation_fixed_point}
\begin{aligned}
    \Delta(\partial_\rho\phi^a)(x)&:= \tfrac{\partial\tilde{\phi}_\epsilon^a(y)}{\partial y^\rho}_{\mid y=x}-\tfrac{\partial{\phi}^a(y)}{\partial y^\rho}_{\mid y=x}
    = \epsilon \delta M^a_b \tfrac{\partial \phi^b(x)}{\partial x^\rho}
    -\epsilon \tfrac{\partial \phi^a(x)}{\partial x^\sigma}\partial_\rho\delta x^\sigma
    -\epsilon \tfrac{\partial^2 \phi^a(x)}{\partial x^\rho \partial x^\nu} \delta x^\nu
    {\overset{!}{=}} \tfrac{\partial}{\partial x^\rho}\Delta\phi^a(x)\,.
\end{aligned}
\end{align}}
Similarly we confirm that variations at the transformed spacetime point and derivatives do not commute since
{\small\begin{align}
\begin{aligned}
    \delta(\partial_\rho\phi^a)(x)&:= \tfrac{\partial\tilde{\phi}_\epsilon^a(y)}{\partial y^\rho}_{\mid y=\tilde{x}}-\tfrac{\partial{\phi}^a(y)}{\partial y^\rho}_{\mid y=x}
    = \epsilon \delta M^a_b \tfrac{\partial \phi^b(x)}{\partial x^\rho}
    -\epsilon \tfrac{\partial \phi^a(x)}{\partial x^\sigma}\partial_\rho\delta x^\sigma
    \neq 
    \epsilon \delta M^a_b \tfrac{\partial \phi^b(x)}{\partial x^\rho}
    =:\tfrac{\partial}{\partial x^\rho}\delta\phi^a(x)\,.
\end{aligned}
\end{align}}
%
\end{proof}

\begin{mytheorem}[Quantum Noether's first theorem \todoinline{finish}]
Consider a quantum field theory with a continuous global symmetry.
If the symmetry is preserved by the quantization procedure and is
non-anomalous, then for every generator $A$ there exists a renormalized
current operator $\hat{j}_A^\mu$ satisfying
{\small\begin{align}
    \partial_\mu\hat{j}_A^\mu=0
\end{align}}
as an operator equation, up to equations of motion and contact terms
in time-ordered correlation functions.

Provided the flux at spatial infinity vanishes, the corresponding
charge
{\small\begin{align}
    \hat{Q}_A[\Sigma]
    :=\int_\Sigma d\Sigma_\mu\,\hat{j}_A^\mu
\end{align}}
is independent of the spacelike hypersurface $\Sigma$. On a
constant-time hypersurface,
{\small\begin{align}
    \hat{Q}_A(t)
    =\int_{\Sigma_t}d^{d-1}\!x\,\hat{j}_A^0(t,\mathbf{x}),
    \qquad
    \frac{d\hat{Q}_A}{dt}=0.
\end{align}}

The conserved charge generates the symmetry on the Hilbert space.
With the convention
{\small\begin{align}
    U(\epsilon)
    =\exp\!\left(-\frac{i}{\hbar}\epsilon^A\hat{Q}_A\right),
    \qquad
    \tilde{\phi}^a
    =U(\epsilon)\hat{\phi}^aU(\epsilon)^{-1},
\end{align}}
the infinitesimal transformation is
{\small\begin{align}
    \Delta_A\hat{\phi}^a(x)
    =-\frac{i}{\hbar}
      [\hat{Q}_A,\hat{\phi}^a(x)].
\end{align}}
In the Heisenberg picture,
{\small\begin{align}
    \frac{d\hat{Q}_A}{dt}
    =
    \frac{\partial\hat{Q}_A}{\partial t}
    +\frac{i}{\hbar}[\hat{H},\hat{Q}_A]
    =0,
\end{align}}
so that, when $\hat{Q}_A$ has no explicit time dependence,
$[\hat{H},\hat{Q}_A]=0$.

At the level of correlation functions these statements are encoded
in the Ward--Takahashi identities. If the symmetry is anomalous, the
current instead satisfies
$\partial_\mu\hat{j}_A^\mu=\hat{\mathcal A}_A\neq0$.
\end{mytheorem}



\begin{mytheorem}[Classical Noether's second theorem \todoinline{finish}]
Let a classical action $S[\phi]$ possess a local gauge symmetry
parametrized by arbitrary smooth functions $\epsilon^\alpha(x)$. For
transformations involving at most first derivatives of the gauge
parameters, write the fixed-point field variation as
{\small\begin{align}
    \Delta_\epsilon\phi^a
    =
    R^a{}_\alpha\,\epsilon^\alpha
    +R^{a\mu}{}_\alpha\,\partial_\mu\epsilon^\alpha .
\end{align}}
Then the Euler--Lagrange expressions
{\small\begin{align}
    \mathcal E_a
    :=
    \frac{\delta S}{\delta\phi^a}
    =
    \frac{\partial\mathcal L}{\partial\phi^a}
    -\partial_\mu
     \frac{\partial\mathcal L}{\partial(\partial_\mu\phi^a)}
\end{align}}
satisfy the off-shell Noether identities
{\small\begin{align}
    \mathcal E_a R^a{}_\alpha
    -\partial_\mu
     \left(\mathcal E_a R^{a\mu}{}_\alpha\right)
    \equiv0 .
\end{align}}
Thus gauge invariance implies that the Euler--Lagrange equations are
not all independent.

The Noether current associated with a local gauge transformation can
generally be written as
{\small\begin{align}
    j^\mu_\epsilon
    =
    W^\mu_\epsilon
    +\partial_\nu k^{[\nu\mu]}_\epsilon,
    \qquad
    W^\mu_\epsilon\approx0 ,
\end{align}}
where $\approx$ denotes equality on shell. Hence its charge reduces to
a boundary term,
{\small\begin{align}
    Q[\epsilon]
    \approx
    \int_{\partial\Sigma}
    dS_{\mu\nu}\,k^{\mu\nu}_\epsilon .
\end{align}}
Gauge transformations that vanish at the boundary therefore carry no
independent physical charge, whereas global or asymptotic gauge
transformations may carry nonzero boundary charges.
\end{mytheorem}



\begin{mytheorem}[Quantum Noether's second theorem \todoinline{finish}]
Consider a quantum gauge theory obtained by gauge fixing a classical
theory with a local symmetry. If the functional measure,
regularization, and renormalization preserve the gauge symmetry, the
quantum theory satisfies Ward--Takahashi or Slavnov--Taylor identities.

Schematically, in a gauge-invariant formulation the renormalized
effective action $\Gamma[\Phi]$ satisfies
{\small\begin{align}
    \frac{\delta\Gamma}{\delta\Phi^a}R^a{}_\alpha
    -
    \partial_\mu\!
    \left(
        \frac{\delta\Gamma}{\delta\Phi^a}
        R^{a\mu}{}_\alpha
    \right)
    =0 .
\end{align}}
In a gauge-fixed formulation the corresponding statement is encoded
by the BRST or Slavnov--Taylor identity
{\small\begin{align}
    \mathcal S(\Gamma)=0 .
\end{align}}

Small gauge transformations are redundancies rather than physical
symmetries. Their generators annihilate physical states,
{\small\begin{align}
    \widehat G[\epsilon]\ket{\Psi_{\mathrm{phys}}}=0 .
\end{align}}
Equivalently, in BRST quantization,
{\small\begin{align}
    \widehat Q_{\mathrm{BRST}}^{\,2}=0,
    \qquad
    \widehat Q_{\mathrm{BRST}}\ket{\Psi_{\mathrm{phys}}}=0,
    \qquad
    \ket{\Psi_{\mathrm{phys}}}
    \sim
    \ket{\Psi_{\mathrm{phys}}}
    +\widehat Q_{\mathrm{BRST}}\ket{\chi}.
\end{align}}
Gauge transformations acting nontrivially at a boundary may instead
generate physical boundary charges.

If the gauge symmetry is anomalous, the quantum identity becomes
{\small\begin{align}
    \mathcal S(\Gamma)
    =\hbar\,\mathcal A\neq0 ,
\end{align}}
and the gauge symmetry cannot be consistently maintained unless the
anomaly is cancelled.
\end{mytheorem}





%--------------------------------------------------------
%=========================================================
\subsection{Energy, Momentum and Angular Momentum}
%=========================================================
%------------------------------------------------------------

Conservation of energy and momentum is a consequence of invariance under spacetime translations $x^\mu\mapsto x^\mu + a^\mu$.
The associated conserved currents are the components of the energy-momentum tensor $T^{\mu\nu}$, which is a symmetric rank-2 tensor that can be derived from the Lagrangian density $\mathcal{L}$ via the Noether procedure.
The conserved charges are the total energy and total mometum 
\begin{align}
E = \int d^3 x\, T^{00}, \quad \mathbf{P} = \int d^3 x\, T^{0i}.
\end{align}

The total energy often coincides with the Hamiltonian $H$ of the system, which generates time evolution, albeit one has to be careful in curved spacetime since several hamiltonian operators can be defined depending on the choice of time coordinate and associated foliation of spacetime, while in the presence of a timelike Killing vector field thee will be a well-defined notion of energy which coincides only with the Hamiltonian defined with respect to the time coordinate adapted to the Killing vector field, thus singling out a preferred foliation of spacetime and a preferred Hamiltonian operator and notion of time evolution.


%==================================================================
\subsubsection[Many-body systems, total momentum, center-of-mass splitting, and finite temperature]{Many-body systems, total momentum, center-of-mass splitting, and finite temperature\todoinline{Fix \& Polish}}
%==================================================================
This note summarizes the main points of the discussion about a many-body Hamiltonian $H$ commuting with the total momentum operator $\mathbf{P}$, the possibility of choosing states with zero total momentum, the explicit center-of-mass transformation, the decomposition $H=H_{\text{cm}}+H_{\text{rel}}$, and the extension to variable particle number and finite temperature.

\subsubsection*{1. What follows from $[H,\mathbf{P}]=0$}
If
{\small\begin{align}
\begin{aligned}
    [H,\mathbf{P}]=0,
\end{aligned}
\end{align}}
then one can choose common eigenstates of $H$ and $\mathbf{P}$:
{\small\begin{align}
\begin{aligned}
    H\ket{E,\mathbf{K}}=E(\mathbf{K})\ket{E,\mathbf{K}},\qquad \mathbf{P}\ket{E,\mathbf{K}}=\mathbf{K}\ket{E,\mathbf{K}}.
\end{aligned}
\end{align}}
However, \emph{this alone does not imply} that the ground state must have $\mathbf{P}=0$.
The statement that the minimum energy occurs at zero total momentum requires more structure, typically Galilean invariance or an explicit decoupling of the center-of-mass motion.
In the standard non-relativistic situation, the energy takes the form
{\small\begin{align}
\begin{aligned}
    E(\mathbf{K})=E_{\mathrm{int}}+\frac{\mathbf{K}^2}{2M},
\end{aligned}
\end{align}}
so the minimum is indeed at $\mathbf{K}=0$.

\subsubsection*{2. Explicit center-of-mass operator and momentum-shifting unitary}
For $N$ particles of masses $m_i$, define
{\small\begin{align}
\begin{aligned}
    M=\sum_{i=1}^N m_i,\qquad \mathbf{P}=\sum_{i=1}^N \mathbf{P}_i,\qquad \mathbf{X}_{\text{cm}}=\frac{1}{M}\sum_{i=1}^N m_i\,\mathbf{X}_i.
\end{aligned}
\end{align}}
Then
{\small\begin{align}
\begin{aligned}
    [X_{\text{cm},a},P_b]=i\delta_{ab}.
\end{aligned}
\end{align}}
The unitary that shifts the total momentum is
{\small\begin{align}
\begin{aligned}
    U(\mathbf{K})=e^{-i\mathbf{K}\cdot \mathbf{X}_{\text{cm}}}.
\end{aligned}
\end{align}}
Using the Baker-Campbell-Hausdorff formula,
{\small\begin{align}
\begin{aligned}
    U(\mathbf{K})^\dagger\mathbf{P}\,U(\mathbf{K})=\mathbf{P}-\mathbf{K}.
\end{aligned}
\end{align}}
Equivalently,
{\small\begin{align}
\begin{aligned}
    \mathbf{P}\,U(\mathbf{K})=U(\mathbf{K})(\mathbf{P}-\mathbf{K}).
\end{aligned}
\end{align}}
Therefore, if
{\small\begin{align}
\begin{aligned}
    \mathbf{P}\ket{\psi_{\mathbf{K}}}=\mathbf{K}\ket{\psi_{\mathbf{K}}},
\end{aligned}
\end{align}}
then the transformed state
{\small\begin{align}
\begin{aligned}
    \ket{\psi_{\mathrm{rest}}}:=U(\mathbf{K})\ket{\psi_{\mathbf{K}}}
\end{aligned}
\end{align}}
obeys
{\small\begin{align}
\begin{aligned}
    \mathbf{P}\ket{\psi_{\mathrm{rest}}}=0.
\end{aligned}
\end{align}}
So $U(\mathbf{K})$ removes the center-of-mass momentum.

\subsubsection*{3. Why this changes the center-of-mass energy}
If the center-of-mass Hamiltonian is
{\small\begin{align}
\begin{aligned}
    H_{\text{cm}}=\frac{\mathbf{P}^2}{2M},
\end{aligned}
\end{align}}
then $\mathbf{X}_{\text{cm}}$ does not commute with $H_{\text{cm}}$:
{\small\begin{align}
\begin{aligned}
    [X_{\text{cm},a},H_{\text{cm}}]=\left[X_{\text{cm},a},\frac{\mathbf{P}^2}{2M}\right]=\frac{iP_a}{M}.
\end{aligned}
\end{align}}
This is precisely why the unitary generated by $\mathbf{X}_{\text{cm}}$ can change the center-of-mass energy.
Indeed,
{\small\begin{align}
\begin{aligned}
    U(\mathbf{K})^\dagger H_{\text{cm}}U(\mathbf{K})=\frac{(\mathbf{P}-\mathbf{K})^2}{2M}.
\end{aligned}
\end{align}}
So if
{\small\begin{align}
\begin{aligned}
    \mathbf{P}\ket{\psi_{\mathbf{K}}}=\mathbf{K}\ket{\psi_{\mathbf{K}}},
\end{aligned}
\end{align}}
then
{\small\begin{align}
\begin{aligned}
    H_{\text{cm}}\,U(\mathbf{K})\ket{\psi_{\mathbf{K}}}=0.
\end{aligned}
\end{align}}
More explicitly, if
{\small\begin{align}
\begin{aligned}
    H_{\text{cm}}\ket{\psi_{\mathbf{K}}}=\frac{\mathbf{K}^2}{2M}\ket{\psi_{\mathbf{K}}},
\end{aligned}
\end{align}}
then after the transformation the center-of-mass energy becomes zero.

\subsubsection*{4. Splitting the free many-body Hamiltonian}
Consider now a free gas of identical particles of mass $m$ in second quantization:
{\small\begin{align}
\begin{aligned}
    H=\sum_{\mathbf{K}}\frac{\mathbf{K}^2}{2m}\,c_{\mathbf{K}}^\dagger c_{\mathbf{K}}.
\end{aligned}
\end{align}}
Define also the total number operator and total momentum operator:
{\small\begin{align}
\begin{aligned}
    \hat N=\sum_{\mathbf{K}} c_{\mathbf{K}}^\dagger c_{\mathbf{K}},\qquad \mathbf{P}=\sum_{\mathbf{K}} \mathbf{K}\,c_{\mathbf{K}}^\dagger c_{\mathbf{K}}.
\end{aligned}
\end{align}}
Write $n_{\mathbf{K}}:=c_{\mathbf{K}}^\dagger c_{\mathbf{K}}$.
Then
{\small\begin{align}
\begin{aligned}
    H=\sum_{\mathbf{K}}\frac{\mathbf{K}^2}{2m}\,n_{\mathbf{K}},\qquad \mathbf{P}=\sum_{\mathbf{K}}\mathbf{K}\,n_{\mathbf{K}}.
\end{aligned}
\end{align}}
Inside a sector of fixed particle number $N$, one may write
{\small\begin{align}
\begin{aligned}
    \mathbf{K}=\left(\mathbf{K}-\frac{\mathbf{P}}{N}\right)+\frac{\mathbf{P}}{N},
\end{aligned}
\end{align}}
and hence
{\small\begin{align}
\begin{aligned}
    \sum_{\mathbf{K}}\mathbf{K}^2 n_{\mathbf{K}}=\sum_{\mathbf{K}}\left(\mathbf{K}-\frac{\mathbf{P}}{N}\right)^2 n_{\mathbf{K}}+2\frac{\mathbf{P}}{N}\cdot \sum_{\mathbf{K}}\left(\mathbf{K}-\frac{\mathbf{P}}{N}\right)n_{\mathbf{K}}+\frac{\mathbf{P}^2}{N^2}\sum_{\mathbf{K}}n_{\mathbf{K}}.
\end{aligned}
\end{align}}
Now
{\small\begin{align}
\begin{aligned}
    \sum_{\mathbf{K}}\left(\mathbf{K}-\frac{\mathbf{P}}{N}\right)n_{\mathbf{K}}=\sum_{\mathbf{K}}\mathbf{K}\,n_{\mathbf{K}}-\frac{\mathbf{P}}{N}\sum_{\mathbf{K}}n_{\mathbf{K}}=\mathbf{P}-\frac{\mathbf{P}}{N}N=0,
\end{aligned}
\end{align}}
and also
{\small\begin{align}
\begin{aligned}
    \sum_{\mathbf{K}}n_{\mathbf{K}}=N.
\end{aligned}
\end{align}}
Therefore,
{\small\begin{align}
\begin{aligned}
    \sum_{\mathbf{K}}\mathbf{K}^2 n_{\mathbf{K}}=\sum_{\mathbf{K}}\left(\mathbf{K}-\frac{\mathbf{P}}{N}\right)^2 n_{\mathbf{K}}+\frac{\mathbf{P}^2}{N},
\end{aligned}
\end{align}}
so that
{\small\begin{align}
\begin{aligned}
    H=\frac{\mathbf{P}^2}{2mN}+\sum_{\mathbf{K}}\frac{\left(\mathbf{K}-\frac{\mathbf{P}}{N}\right)^2}{2m}\,n_{\mathbf{K}}.
\end{aligned}
\end{align}}
This is the desired decomposition
{\small\begin{align}
\begin{aligned}
    H=H_{\text{cm}}+H_{\text{rel}},
\end{aligned}
\end{align}}
with
{\small\begin{align}
\begin{aligned}
    H_{\text{cm}}=\frac{\mathbf{P}^2}{2M},\qquad M=Nm,
\end{aligned}
\end{align}}
and
{\small\begin{align}
\begin{aligned}
    H_{\text{rel}}=\sum_{\mathbf{K}}\frac{\left(\mathbf{K}-\frac{\mathbf{P}}{N}\right)^2}{2m}\,c_{\mathbf{K}}^\dagger c_{\mathbf{K}}.
\end{aligned}
\end{align}}
Thus, at fixed $N$,
{\small\begin{align}
\begin{aligned}
    E_{N,\alpha}(\mathbf{K})=E^{\text{rel}}_{N,\alpha}+\frac{\mathbf{K}^2}{2mN},
\end{aligned}
\end{align}}
so the minimum in total momentum is attained at $\mathbf{K}=0$.

\subsubsection*{4.bis. Extension to interacting systems}
More generally the argument goes through for any Hamiltonian that can be written as
{\small\begin{align}
\begin{aligned}
    H=H_{\text{cm}}+H_{\text{rel}}=\frac{\mathbf{P}^2}{2M}+H_{\text{rel}},
\end{aligned}
\end{align}}
with $[H_{\text{rel}},\mathbf{P}]=0$.
This includes a large class of interacting systems, such as those with two-body interactions that depend only on the relative coordinates, e.g.
{\small\begin{align}
\begin{aligned}
    H=\sum_{i=1}^N \frac{\mathbf{P}_i^2}{2m}+\sum_{i<j}V(\mathbf{X}_i-\mathbf{X}_j).
\end{aligned}\end{align}}
In this case, the energy still takes the form
{\small\begin{align}
\begin{aligned}
    E(\mathbf{K})=E_{\text{rel}}+\frac{\mathbf{K}^2}{2M},
\end{aligned}\end{align}}
so the minimum is at $\mathbf{K}=0$.
The key point is that the center-of-mass coordinate $\mathbf{X}_{\text{cm}}$ still generates momentum shifts, and the center-of-mass Hamiltonian still has the form $H_{\text{cm}}=\mathbf{P}^2/(2M)$, so the same argument applies.


\subsubsection*{4.tris. Extension to interacting systems REMASTERED}
Summarizing, the general Hamiltonian will receive single particle and interaction contributions, single particle contributions may also include the interaction with an external potential, and the interaction contributions may include two-body, three-body, and higher-body interactions, so the most general form of the Hamiltonian is
\begin{align}
   H= \sum_k \tfrac{k^2}{2m}c_k^\dagger c_k + \sum_{ab}h_{ab}c_a^\dagger c_b + \sum_{abcd}V_{abcd}c_a^\dagger c_b^\dagger c_c c_d + \sum_{abcdef}W_{abcdef}c_a^\dagger c_b^\dagger c_c^\dagger c_d c_e c_f + \cdots,
\end{align}
We can always split the single-particle kinetic term as
\begin{align}
   \sum_k \tfrac{k^2}{2m}c_k^\dagger c_k = \frac{\mathbf{P}^2}{2mN} + \sum_k \tfrac{(k-\mathbf{P}/N)^2}{2m}c_k^\dagger c_k
\end{align}
and the remaining single-particle and interaction terms can be rearranged into a relative Hamiltonian $H_{\text{rel}}$
\begin{align}
    H= \frac{\mathbf{P}^2}{2mN} + \sum_k \tfrac{(k-\mathbf{P}/N)^2}{2m}c_k^\dagger c_k + \sum_{ab}h_{ab}c_a^\dagger c_b + \sum_{abcd}V_{abcd}c_a^\dagger c_b^\dagger c_c c_d + \sum_{abcdef}W_{abcdef}c_a^\dagger c_b^\dagger c_c^\dagger c_d c_e c_f + \cdots =: \frac{\mathbf{P}^2}{2mN} + H_{\text{rel}}.
\end{align}
If the assumption $[H,\mathbf{P}]=0$ holds, then also $[H_{\text{rel}},\mathbf{P}]=0$.
In particular, we can classify states by their total momentum $\mathbf{K}$, and the energy takes the form
\begin{align}
    E_{N,\alpha}(\mathbf{K})=E^{\text{rel}}_{N,\alpha}+\frac{\mathbf{K}^2}{2mN},
\end{align}
so the minimum is at $\mathbf{K}=0$.

\subsubsection*{4.quater. Aside on the relativistic case}\todotag{Very not sure...}
In a relativistic theory, the energy-momentum relation is $E^2=\mathbf{P}^2+M^2$, so the energy takes the form
{\small\begin{align}
\begin{aligned}
    E(\mathbf{K})=\sqrt{\mathbf{K}^2+M^2},
\end{aligned}\end{align}}
and the minimum is still at $\mathbf{K}=0$.
The key point is that the center-of-mass coordinate still generates momentum shifts, and the center-of-mass Hamiltonian still has the form $H_{\text{cm}}=\sqrt{\mathbf{P}^2+M^2}$, so the same argument applies.
The main difference with the non-relativistic case is that the center-of-mass energy is not quadratic in $\mathbf{P}$, but the minimum is still at zero total momentum.


\subsubsection*{5. First-quantized version}
For $N$ identical particles,
{\small\begin{align}
\begin{aligned}
    H=\sum_{i=1}^N \frac{\mathbf{P}_i^2}{2m},\qquad \mathbf{P}=\sum_{i=1}^N \mathbf{P}_i.
\end{aligned}
\end{align}}
Define the relative momenta
{\small\begin{align}
\begin{aligned}
    \boldsymbol{\pi}_i:=\mathbf{P}_i-\frac{\mathbf{P}}{N}.
\end{aligned}
\end{align}}
Then
{\small\begin{align}
\begin{aligned}
    \sum_{i=1}^N \boldsymbol{\pi}_i=0,
\end{aligned}
\end{align}}
and
{\small\begin{align}
\begin{aligned}
    \sum_{i=1}^N \mathbf{P}_i^2=\sum_{i=1}^N\left(\boldsymbol{\pi}_i+\frac{\mathbf{P}}{N}\right)^2=\sum_{i=1}^N \boldsymbol{\pi}_i^2+\frac{\mathbf{P}^2}{N}.
\end{aligned}
\end{align}}
Therefore,
{\small\begin{align}
\begin{aligned}
    H=\frac{\mathbf{P}^2}{2mN}+\sum_{i=1}^N\frac{\boldsymbol{\pi}_i^2}{2m}.
\end{aligned}
\end{align}}
This is the same center-of-mass/relative splitting in a more elementary form.
For $N=2$, with
{\small\begin{align}
\begin{aligned}
    \mathbf{K}=\mathbf{K}_1+\mathbf{K}_2,\qquad \mathbf{q}=\frac{\mathbf{K}_1-\mathbf{K}_2}{2},
\end{aligned}
\end{align}}
one has
{\small\begin{align}
\begin{aligned}
    \mathbf{K}_1=\frac{\mathbf{K}}{2}+\mathbf{q},\qquad \mathbf{K}_2=\frac{\mathbf{K}}{2}-\mathbf{q},
\end{aligned}
\end{align}}
and therefore
{\small\begin{align}
\begin{aligned}
    \frac{\mathbf{K}_1^2}{2m}+\frac{\mathbf{K}_2^2}{2m}=\frac{\mathbf{K}^2}{4m}+\frac{\mathbf{q}^2}{m}.
\end{aligned}
\end{align}}
Again, the first term is the center-of-mass kinetic energy and the second is the relative one.

\subsubsection*{6. Variable particle number and Fock-space decomposition}
If the Hamiltonian commutes both with the total momentum and with the particle number,
{\small\begin{align}
\begin{aligned}
    [H,\mathbf{P}]=0,\qquad [H,\hat N]=0,
\end{aligned}
\end{align}}
then the Fock space decomposes into sectors of fixed $N$,
{\small\begin{align}
\begin{aligned}
    \mathcal F=\bigoplus_{N=0}^{\infty}\mathcal H_N,
\end{aligned}
\end{align}}
and one may choose common eigenstates of $H$, $\mathbf{P}$, and $\hat N$:
{\small\begin{align}
\begin{aligned}
    H\ket{N,\alpha,\mathbf{K}}=E_{N,\alpha}(\mathbf{K})\ket{N,\alpha,\mathbf{K}},
\end{aligned}
\end{align}}
{\small\begin{align}
\begin{aligned}
    \mathbf{P}\ket{N,\alpha,\mathbf{K}}=\mathbf{K}\ket{N,\alpha,\mathbf{K}},\qquad \hat N\ket{N,\alpha,\mathbf{K}}=N\ket{N,\alpha,\mathbf{K}}.
\end{aligned}
\end{align}}
Inside each sector with $N\geq 1$, the center-of-mass splitting still holds:
{\small\begin{align}
\begin{aligned}
    E_{N,\alpha}(\mathbf{K})=E^{\text{rel}}_{N,\alpha}+\frac{\mathbf{K}^2}{2mN}.
\end{aligned}
\end{align}}
Hence, for each fixed $N$, the minimum in total momentum occurs at $\mathbf{K}=0$.
However, if one minimizes the bare Hamiltonian $H$ over the entire Fock space, the result can be trivial.
For the free gas,
{\small\begin{align}
\begin{aligned}
    H=\sum_{\mathbf{K}}\frac{\mathbf{K}^2}{2m}c_{\mathbf{K}}^\dagger c_{\mathbf{K}}\geq 0,
\end{aligned}
\end{align}}
and the vacuum has zero energy, so the global ground state of $H$ is simply the vacuum:
{\small\begin{align}
\begin{aligned}
    \ket{0},\qquad \hat N\ket{0}=0,\qquad \mathbf{P}\ket{0}=0.
\end{aligned}
\end{align}}
So the physically interesting many-body problem is usually either the fixed-$N$ problem or the grand-canonical one.
A careful operator-level way to write the center-of-mass piece on the whole Fock space is sector by sector:
{\small\begin{align}
\begin{aligned}
    H_{\text{cm}}=\sum_{N\geq 1}\Pi_N\frac{\mathbf{P}^2}{2mN}\Pi_N,
\end{aligned}
\end{align}}
with $H_{\text{cm}}\ket{0}=0$.
Here $\Pi_N$ denotes the projector onto the $N$-particle sector.

\subsubsection*{7. Finite temperature and the grand-canonical point of view}
At finite temperature, the main object is no longer a single ground state but a thermal density matrix.
If $N$ is fixed, one uses the canonical ensemble,
{\small\begin{align}
\begin{aligned}
    \rho_N=\frac{e^{-\beta H^{(N)}}}{Z_N}.
\end{aligned}
\end{align}}
If $N$ is allowed to fluctuate, the natural object is the grand-canonical ensemble,
{\small\begin{align}
\begin{aligned}
    \rho=\frac{e^{-\beta(H-\mu \hat N)}}{\Xi}.
\end{aligned}
\end{align}}
The operator that matters is then
{\small\begin{align}
\begin{aligned}
    K:=H-\mu \hat N.
\end{aligned}
\end{align}}
Inside the sector with fixed $N$,
{\small\begin{align}
\begin{aligned}
    K^{(N)}=H^{(N)}-\mu N,
\end{aligned}
\end{align}}
so the eigenvalues take the form
{\small\begin{align}
\begin{aligned}
    E^{(K)}_{N,\alpha}(\mathbf{K})=E^{\text{rel}}_{N,\alpha}+\frac{\mathbf{K}^2}{2mN}-\mu N.
\end{aligned}
\end{align}}
Again, for every fixed $N$, the minimum in total momentum occurs at
{\small\begin{align}
\begin{aligned}
    \mathbf{K}=0.
\end{aligned}
\end{align}}
What changes is which particle-number sector is thermodynamically favored: this is determined by the competition between the internal energy and the term $-\mu N$.
In equilibrium, unless a macroscopic flow is imposed, one typically has
{\small\begin{align}
\begin{aligned}
    \langle \mathbf{P}\rangle=0,
\end{aligned}
\end{align}}
not because each pure state in the thermal mixture has zero momentum, but because the thermal weights are symmetric and the lowest-energy contribution in each fixed-$N$ sector lies at zero total momentum.

\subsubsection*{8. Final summary}
The main logical points are the following.
\begin{itemize}
    \item From $[H,\mathbf{P}]=0$ alone, one may choose common eigenstates of $H$ and $\mathbf{P}$, but one cannot conclude by pure algebra that the ground state must have zero momentum.
    \item In Galilean-invariant non-relativistic systems, the center-of-mass coordinate
    {\small\begin{align}
    \begin{aligned}
        \mathbf{X}_{\text{cm}}=\frac{1}{M}\sum_i m_i\,\mathbf{X}_i
    \end{aligned}
    \end{align}}
    generates momentum shifts via
    {\small\begin{align}
    \begin{aligned}
        U(\mathbf{K})=e^{-i\mathbf{K}\cdot \mathbf{X}_{\text{cm}}},
    \end{aligned}
    \end{align}}
    and this maps a state of total momentum $\mathbf{K}$ to one of total momentum zero.
    \item When the Hamiltonian splits as
    {\small\begin{align}
    \begin{aligned}
        H=H_{\text{cm}}+H_{\text{rel}}=\frac{\mathbf{P}^2}{2M}+H_{\text{rel}},
    \end{aligned}
    \end{align}}
    the energy is
    {\small\begin{align}
    \begin{aligned}
        E(\mathbf{K})=E_{\text{rel}}+\frac{\mathbf{K}^2}{2M},
    \end{aligned}
    \end{align}}
    so the minimum is at $\mathbf{K}=0$.
    \item In second quantization, the splitting holds cleanly inside each sector of fixed particle number $N$:
    {\small\begin{align}
    \begin{aligned}
        H=\frac{\mathbf{P}^2}{2mN}+H_{\text{rel}}.
    \end{aligned}
    \end{align}}
    \item If the particle number is allowed to vary, the global ground state of the free Hamiltonian $H$ is just the vacuum.
    The physically relevant finite-density problem is either fixed-$N$ or grand canonical, where the relevant operator is $H-\mu\hat N$.
    \item At finite temperature, the correct object is the thermal density matrix, not a single ground state.
    Still, in each fixed-$N$ sector, the energetically preferred total momentum remains zero.
\end{itemize}
\end{document}
